\section{An Euler-characteristic bound on the resolution}\label{sec:euler}

We now prove the geometric bound which permits a complete universal
calculation from the genus-zero data. Continue with the resolution
$Y_+$ and the integral nef divisors of Section~\ref{sec:geometry}.

\begin{theorem}\label{thm:euler}
Every pure one-dimensional Cohen--Macaulay subscheme $Z\subset Y_+$
with
\[
 H_4[Z]=0,\qquad (H_2[Z],H_3[Z])\in\{(1,0),(0,1),(1,1)\}
\]
satisfies $\chi(\mathcal O_Z)\ge1$. No connectedness assumption or
bound on the multiplicities of exceptional components is required.
\end{theorem}

\subsection{The contracted surfaces}
The toric divisor $D_0$ restricts to $H_4$. Its convex integral support
function defines a semiample contraction. The minimal cones whose
orbits have positive-dimensional fibers are
\begin{equation}\label{eq:contracted}
 (8),\ (9),\ (10),\ (2,4),\ (6,7).
\end{equation}
This is a finite face calculation from the fan in Appendix~\ref{app:data}:
form $P_0=\{m:\langle m,v_i\rangle+\delta_{i0}\ge0\}$.
For a cone $\sigma$, impose its equalities to obtain $F_\sigma$.
The image orbit cone has dimension $4-\dim F_\sigma$; the orbit fiber
is positive-dimensional exactly when this exceeds $\dim\sigma$.
Taking minimal cones gives \eqref{eq:contracted}.

The divisor $D_{10}$ misses the general hypersurface, since its dual
Newton face is a vertex with nonzero coefficient. Thus every
$H_4$-degree-zero curve is supported on
\[
 E\cup F\cup C_1\cup C_2,
 \qquad E=D_9|_{Y_+},\quad F=D_8|_{Y_+},
\]
where $C_1,C_2$ come from the orbit closures $67,24$.
There are no cones containing $678$ or $248$, and the Newton faces at
$679,249$ are singletons. Hence the $C_i$ miss both surfaces.
The absence of a cone containing $2467$ makes them disjoint.

We need the actual surfaces, rather than only their intersection forms:
\begin{equation}\label{eq:surfaces}
 E\simeq\Bl_{p_1,p_2}\PP^2,\qquad F\simeq\FF_1.
\end{equation}
For $E$, the anticanonical face is the primitive edge with endpoints
$(0,1,-1,-1)$ and $(1,-1,-1,-1)$. Its binomial fixes a nonzero torus
parameter. Since $D_{10}$ is absent, rays $2,7$ are also excluded
on $E$. The remaining rays $3,4,5,6,8$ project respectively to
\[
 (1,0),\quad(0,1),\quad(-1,0),\quad(0,-1),\quad(1,1),
\]
with cones
$38,84,45,56,63$, the fan of the plane blown up at two distinct
torus fixed points.
For $F$, the star away from $D_{10}$ has cones
$137,147,379,479$. Quotient by the fiber ray $7$ gives base rays
$(1,0),(0,1),(-1,-1),(-1,0)$, the fan of $\FF_1$.
The Newton vertex $(1,-1,-1,-1)$ has exponents
$(0,0,0,1,0,0)$ on rays $1,3,4,7,9,10$; all other vertices have
fiber exponent zero. The equation is therefore monic linear in the
fiber coordinate, and its zero locus is a graph over $\FF_1$.
This proves \eqref{eq:surfaces}.

Write $h$ for the pullback of a plane line and $e_1,e_2$ for the
exceptional curves of $E$. The two surfaces intersect transversely in
$s=h-e_1-e_2$, which is the negative section on $F$. The curves
$e_i$ meet $s$ transversely at distinct points $q_i$. All four curves
$e_1,e_2,C_1,C_2$ are pairwise disjoint $(-1,-1)$-curves.
For $e_i$ the normal sequence has subbundle and quotient
$\mathcal O(-1)$, since $N_{E/Y_+}=K_E$; it splits by
$H^1(\mathcal O_{\PP^1})=0$. For $C_i$ this is the local node model.
The integral class identities are
\begin{equation}\label{eq:surface-classes}
 \begin{gathered}
 s=b_2,\quad f=b_3,\quad e_1=b_1+b_5+b_6,\quad e_2=b_6,\\
 h=b_1+b_2+b_5+2b_6,\qquad [C_1]=b_5,\quad[C_2]=b_1.
 \end{gathered}
\end{equation}
Thus the bidegree $(H_2,H_3)$ of a curve in $E$ is $(h\cdot C,0)$,
and that of $bs+af$ on $F$ is $(b,a)$. The only irreducible curves
of bidegree zero in the contracted locus are the four just listed.

\subsection{Multiple exceptional components}
\begin{lemma}\label{lem:multiple}
Let $M$ be a pure Cohen--Macaulay curve supported with multiplicity
$m$ on a smooth rational $(-1,-1)$-curve $C$ in a smooth threefold.
Then $\chi(\mathcal O_M)\ge2m-1$. If a reduced curve $B$ near
$p\in C$ lies in a smooth surface transverse to $C$, then
$\operatorname{length}_p(M\cap B)\le m$.
\end{lemma}
\begin{proof}
The finite conormal filtration has quotients
$G_k=I_C^k\mathcal O_M/I_C^{k+1}\mathcal O_M$, each a quotient of
$\mathcal O_C(k)^{\oplus(k+1)}$. If its generic rank is $r_k$,
its torsion-free quotient has degree at least $kr_k$, because its
twist by $-k$ is globally generated. Torsion adds nonnegative length,
so $\chi(G_k)\ge(k+1)r_k$. Since $G_0=\mathcal O_C$ and
$\sum r_k=m$, summation gives $\chi(\mathcal O_M)\ge1+2(m-1)$.

If $t$ is an equation of the transverse surface, the completed local
ring of $M$ is finite over $\CC[[t]]$, and is torsion-free because
$M$ is Cohen--Macaulay and $t$ is a parameter. It is free of rank
$m$, so its quotient by $t$ has length $m$. Since $t$ vanishes on
$B$, the asserted intersection length is at most $m$. The argument
applies to several branches together whenever they lie in the same
transverse smooth surface.
\end{proof}

\subsection{Positive-degree components and attachments}
\begin{proof}[Proof of Theorem~\ref{thm:euler}]
A pure curve has no embedded points. Isolated primary decomposition
therefore expresses $Z$ as
$B\cup M_1\cup\cdots\cup M_k$, where each $M_i$ is a pure multiple
structure on one of the four disjoint exceptional curves. Every
positive-bidegree component has generic multiplicity one by nefness
and the stated bidegrees. Its primary ideal is the contraction of its
generic localization, of length one, so that primary component is
reduced. Their union is $B$.

The possibilities for $B$ are
\[
\begin{array}{c|l}
 (1,0)&\text{a strict transform }L\text{ of a plane line, including }s,\\
 (0,1)&\text{a ruling fiber }f\text{ of }F,\\
 (1,1)&\text{an irreducible member of }|s+f|\text{ on }F,
             \text{ or }L\cup f.
\end{array}
\]
Indeed a degree-one curve under $E\to\PP^2$ is a line, and the
effective cone of $\FF_1$ is generated by $s,f$. All the listed
irreducible curves are smooth rational. An irreducible member of
$|s+f|$ misses $s$, and thus misses $E$ and the $e_i$.

A line of class $h$ misses the $e_i$ and meets $s$ once away from
$q_1,q_2$. A line of class $h-e_i$ meets $e_i$ once away from $q_i$
and misses $s$. The line $s$ meets both $e_i$ at the $q_i$.
A ruling fiber meets $s$ once, hence at most one $e_i$.
All these intersections are transverse. Thus $L\cup f$ is either
disjoint, with $\chi(\mathcal O_B)=2$, or has one transverse node,
with $\chi(\mathcal O_B)=1$.

Only one configuration permits two distinct attachments to an
exceptional curve: $L$ has class $h-e_i$ and $f$ passes through
$q_i$. Then $L$ and $f$ are disjoint and $\chi(\mathcal O_B)=2$;
at most one $e_i$ has two attachments. If instead $L=s$ and $f$
passes through $q_i$, both branches lie in the smooth surface $F$
transverse to $e_i$. Lemma~\ref{lem:multiple} bounds their combined
intersection length by $m_i$, rather than $2m_i$.

The union exact sequence now gives
\[
 \chi(\mathcal O_Z)=\chi(\mathcal O_B)+
 \sum_i\bigl(\chi(\mathcal O_{M_i})-
                    \operatorname{length}(B\cap M_i)\bigr).
\]
A multiple component with zero or one attachment contributes at
least zero. A component with two distinct attachments contributes
at least $-1$; there is at most one, and it occurs only when
$\chi(\mathcal O_B)=2$. The result follows in every case.
\end{proof}

For a stable pair $\mathcal O_{Y_+}\to\mathcal F$, the image is
$\mathcal O_Z$ for a pure curve $Z$, and the cokernel is
zero-dimensional. The theorem therefore implies
$\chi(\mathcal F)\ge1$. Every stable-pair series in the three
bidegrees starts in positive Euler characteristic, including those
with arbitrarily large exceptional multiplicities.
